%% ============================================================
%%  SHAMEL — نظام الحضور الذكي
%%  Graduation Thesis — Main Template
%%  Compiler: XeLaTeX  (Menu → Settings → Compiler → XeLaTeX)
%% ============================================================

\documentclass[12pt, a4paper]{report}

%% ─── Core Encoding & Fonts ────────────────────────────────
\usepackage{fontspec}
\usepackage{polyglossia}
\setmainlanguage[numerals=maghrib]{arabic}
\setotherlanguage{english}

%% FreeSerif/FreeSans/FreeMono متاحة في Overleaf وتدعم العربية
\setmainfont{FreeSerif}[Script=Arabic, Scale=1.1]
\setsansfont{FreeSans}[Script=Arabic]
\setmonofont{FreeMono}[Scale=0.88]

\newfontfamily\arabicfont{FreeSerif}[Script=Arabic, Scale=1.1]
\newfontfamily\arabicfontsf{FreeSans}[Script=Arabic]
\newfontfamily\arabicfonttt{FreeMono}[Scale=0.88]
\newfontfamily\englishfont{Latin Modern Roman}

%% ─── Page Geometry ────────────────────────────────────────
\usepackage[
  a4paper,
  top=2.8cm, bottom=2.5cm,
  left=2.5cm, right=2.8cm,
  headheight=15pt
]{geometry}

%% ─── Line Spacing ─────────────────────────────────────────
\usepackage{setspace}
\onehalfspacing

%% ─── Graphics & Figures ───────────────────────────────────
\usepackage{graphicx}
\usepackage{float}
\usepackage{caption}
\captionsetup{font=small, labelfont=bf, justification=centering}
\usepackage{subcaption}

%% ─── TikZ (with externalization to beat timeout) ─────────
\usepackage{tikz}
\usepackage{pgf}
\usetikzlibrary{
  shapes.geometric, arrows.meta, positioning,
  fit, backgrounds, calc, decorations.pathreplacing,
  matrix, chains, shadows
}
\usepackage{pgfplots}
\pgfplotsset{compat=1.18}

%% TikZ externalize — يخزن كل مخطط كـ PDF منفصل
%% يُسرِّع إعادة التجميع بشكل كبير
\usetikzlibrary{external}
\tikzexternalize[prefix=tikz-cache/]

%% ─── Colors ───────────────────────────────────────────────
\usepackage{xcolor}
\definecolor{shamelnavy}{HTML}{0B2545}
\definecolor{shamelgold}{HTML}{C9A84C}
\definecolor{shamelemr}{HTML}{15803D}
\definecolor{codebg}{HTML}{F8F9FA}
\definecolor{codeframe}{HTML}{DEE2E6}
\definecolor{codekeyword}{HTML}{0B2545}
\definecolor{codestring}{HTML}{15803D}
\definecolor{codecomment}{HTML}{6B7280}
\definecolor{codenumber}{HTML}{B45309}

%% ─── Tables ───────────────────────────────────────────────
\usepackage{booktabs}
\usepackage{longtable}
\usepackage{tabularx}
\usepackage{multirow}
\usepackage{array}

%% ─── Code Listings ────────────────────────────────────────
\usepackage{listings}

\lstdefinestyle{pythonstyle}{
  backgroundcolor=\color{codebg},
  frame=single, framerule=0.5pt,
  rulecolor=\color{codeframe},
  basicstyle=\ttfamily\small,
  keywordstyle=\color{codekeyword}\bfseries,
  stringstyle=\color{codestring},
  commentstyle=\color{codecomment}\itshape,
  numberstyle=\tiny\color{codenumber},
  numbers=left, numbersep=8pt,
  tabsize=4, showstringspaces=false,
  breaklines=true, language=Python,
  morekeywords={self,True,False,None,async,await},
  xleftmargin=15pt, aboveskip=8pt, belowskip=8pt,
}

\lstdefinestyle{jsstyle}{
  backgroundcolor=\color{codebg},
  frame=single, framerule=0.5pt,
  rulecolor=\color{codeframe},
  basicstyle=\ttfamily\small,
  keywordstyle=\color{codekeyword}\bfseries,
  stringstyle=\color{codestring},
  commentstyle=\color{codecomment}\itshape,
  numbers=left, numbersep=8pt,
  tabsize=2, showstringspaces=false,
  breaklines=true, language=JavaScript,
  xleftmargin=15pt, aboveskip=8pt, belowskip=8pt,
}

\lstdefinestyle{bashstyle}{
  backgroundcolor=\color{codebg},
  frame=single, framerule=0.5pt,
  rulecolor=\color{codeframe},
  basicstyle=\ttfamily\small,
  keywordstyle=\color{codekeyword}\bfseries,
  commentstyle=\color{codecomment}\itshape,
  numbers=left, numbersep=8pt,
  showstringspaces=false,
  breaklines=true, language=bash,
  xleftmargin=15pt, aboveskip=8pt, belowskip=8pt,
}
\lstset{style=pythonstyle}

%% ─── Math ─────────────────────────────────────────────────
\usepackage{amsmath, amssymb, mathtools}

%% ─── References & Links ───────────────────────────────────
\usepackage{hyperref}
\hypersetup{
  colorlinks=true,
  linkcolor=shamelnavy,
  citecolor=shamelemr,
  urlcolor=shamelnavy,
  pdftitle={SHAMEL - نظام الحضور الذكي},
  bookmarksnumbered=true,
}
\usepackage[backend=biber, style=ieee, sorting=none]{biblatex}
\addbibresource{references.bib}

%% ─── Headers & Footers ────────────────────────────────────
\usepackage{fancyhdr}
\pagestyle{fancy}
\fancyhf{}
\fancyhead[L]{\small\leftmark}
\fancyhead[R]{\small SHAMEL}
\fancyfoot[C]{\thepage}
\renewcommand{\headrulewidth}{0.4pt}

%% ─── Misc ─────────────────────────────────────────────────
\usepackage{enumitem}
\usepackage{csquotes}
\usepackage{tcolorbox}
\tcbuselibrary{breakable, skins}

\newtcolorbox{infobox}[1][]{
  colback=codebg, colframe=shamelnavy,
  fonttitle=\bfseries, title=#1,
  arc=4pt, boxrule=0.6pt
}

%% ─── Chapter Style ────────────────────────────────────────
\usepackage{titlesec}
\titleformat{\chapter}[display]
  {\normalfont\huge\bfseries\color{shamelnavy}}
  {\chaptertitlename\ \thechapter}{16pt}{\Huge}
\titlespacing*{\chapter}{0pt}{20pt}{15pt}

%% ═══════════════════════════════════════════════════════════
\begin{document}
%% ═══════════════════════════════════════════════════════════

%% ─── Cover Page ───────────────────────────────────────────
\begin{titlepage}
\centering
\vspace*{0.8cm}
{\large\bfseries جامعة السودان للعلوم والتكنولوجيا}\\[0.2cm]
{\large كلية علوم الحاسوب والرياضيات}\\[0.2cm]
{\large قسم علم الحاسوب}\\[1.5cm]

\begin{tikzpicture}[remember picture]
  \draw[fill=shamelnavy, draw=shamelgold, line width=3pt] circle (1.8cm);
  \node[white, font=\fontsize{36}{44}\selectfont\bfseries] at (0,0) {ش};
  \draw[shamelgold, line width=1.5pt] circle (2.2cm);
\end{tikzpicture}\\[1.2cm]

{\fontsize{28}{34}\selectfont\bfseries\color{shamelnavy} شامل}\\[0.3cm]
{\Large\color{shamelgold} نظام الحضور الذكي المتكامل}\\[0.2cm]
{\large\color{gray} SHAMEL — Smart Holistic Attendance Management}\\[1.5cm]

{\large بحث مقدم لاستيفاء متطلبات الحصول على درجة البكالوريوس}\\[0.2cm]
{\large في علم الحاسوب}\\[1.5cm]

\begin{tabular}{rl}
  \textbf{إعداد:} & أحمد \\[0.3cm]
  \textbf{إشراف:} & أ.د. \\[0.3cm]
  \textbf{العام الدراسي:} & 2025 -- 2026 م
\end{tabular}
\vfill
{\small\color{gray} نظام شامل — \texttt{https://shamel.sd}}
\end{titlepage}

%% ─── Front Matter ─────────────────────────────────────────
\pagenumbering{roman}

\chapter*{الملخص}
\addcontentsline{toc}{chapter}{الملخص}
يُقدِّم هذا البحث نظام \textbf{شامل} (SHAMEL)، منظومة متكاملة لإدارة الحضور الأكاديمي الذكي باستخدام تقنيات التعرف على الوجه وتحليل البيانات الفوري. يعتمد النظام على Django 6.0.1 مع ASGI وWebSocket Channels، ويُدمج مكتبات DeepFace وdlib للتعرف على الوجه بدقة عالية. يوفر تطبيق ويب تقدمي (PWA) يدعم العمل دون إنترنت عبر تزامن SQLite-PostgreSQL، ولوحة تحكم تفاعلية لأربعة أدوار: المسؤول والمدرس والطالب وموظف الأمن. تم نشر النظام على VPS باستخدام Docker Compose وNginx وSSL على نطاق \texttt{shamel.sd}.

\bigskip\noindent\textbf{الكلمات المفتاحية:} التعرف على الوجه، الحضور الذكي، Django Channels، PWA، PostgreSQL، Docker.

\bigskip\noindent\textbf{Abstract:} SHAMEL is a comprehensive smart attendance management system leveraging facial recognition and real-time analytics, built on Django 6.0.1 with ASGI/WebSocket channels, integrating DeepFace/dlib AI libraries, PWA offline support, and Docker production deployment at shamel.sd.

\chapter*{الإهداء}
\addcontentsline{toc}{chapter}{الإهداء}
\begin{center}\vspace*{2cm}
  {\Large إلى أسرتي الكريمة وأساتذتي الأجلاء}\\[0.5cm]
  {\large إلى كل من ساهم في إنجاز هذا العمل}
\end{center}

\chapter*{شكر وتقدير}
\addcontentsline{toc}{chapter}{شكر وتقدير}
الحمد لله على توفيقه لإتمام هذا العمل. والشكر الجزيل للمشرف الكريم على توجيهاته، ولكلية علوم الحاسوب على ما وفرته من بيئة بحثية داعمة.

\tableofcontents
\listoffigures
\listoftables

%% ─── Chapters ─────────────────────────────────────────────
\cleardoublepage
\pagenumbering{arabic}

\include{chapters/introduction}
\include{chapters/literature_review}
\include{chapters/system_design}
\include{chapters/implementation}
\include{chapters/conclusion}

%% ─── Bibliography ─────────────────────────────────────────
\printbibliography[heading=bibintoc, title={المراجع والمصادر}]

%% ─── Appendix ─────────────────────────────────────────────
\appendix
\include{chapters/appendix_a}

\end{document}
